\problemname{Arithmetic Adaptation}

\illustration{0.3}{regnestykke.jpg}{
  Postcard showing a Norwegian by solving Sample~3.
  Coloured black and white photo,
  National Library of Norway.
}

\newcommand{\maxa}{999}
\newcommand{\maxb}{999}
\newcommand{\maxs}{999}

Practice does make perfect!
You have finally achieved proficiency at the task of adding two small nonzero integers~$a$ and~$b$ to compute their sum~$a+b$.
Before you move on to studying four-digit numbers,
you want to achieve mastery by also understanding the inverse problem: 
given integer~$s$, determine nonzero $a$ and $b$ such that $a+b=s$.
None of the numbers may use more than $3$~digits.


\begin{Input}
    The input consists of:
    \begin{itemize}
        \item One line with an integer $s$ such that $-\maxs\leq s\leq \maxs$.
    \end{itemize}
\end{Input}

\begin{Output}
  Output two integers
  $a$ (with $-\maxa\leq a\leq \maxa$ and $a\neq 0$) 
  and 
  $b$ (with $-\maxb\leq b\leq \maxb$ and $b\neq 0$) 
  such that $a+b=s$.
  If there is more than one valid solution, you may output any one of them.
\end{Output}

\subsection*{Explanation of Sample Input 1}

On input ``\texttt{10}'', the output ``\texttt{3 7}'' is correct because $3+7=10$.
Note that many other outputs would also be correct, such as ``\texttt{2 8}'', ``\texttt{11 -1}'', or even ``\texttt{-849 859}''.
On the other hand, the answer ``\texttt{4 7}'' would be \emph{wrong} (because $4+7\neq 10$), 
and so would ``\texttt{10 0}'' (because both $a$ and $b$ must be nonzero) 
and ``\texttt{1000 -990}'' (because both $a$ and $b$ must have at most three digits.)
